\begin{theorem}[K3 collapse of the nonzero toric fibers]
\label{thm:toric-fiber-K3}
Let $p\geq5$ be prime and $a\in\mathbb F_p^\times$.  Retain the
notation $\mu_p(a)$ from Proposition~\ref{prop:toric-mellin-square}, and put
\[
 \varepsilon_a=\mathbf 1_{a=1},\qquad
 \delta_a=\mathbf 1_{a^2-34a+1=0}.
\]
There is a K3 surface $\widetilde X_a/\mathbb F_p$ such that, if
\[
 \tau_p(a)=\operatorname {Tr}\!\left(
  \operatorname {Frob}_p\mid
  H^2_{\mathrm{et}}(\widetilde X_{a,\overline{\mathbb F}_p},
                    \mathbb Q_\ell)
 \right)
 \qquad(\ell\ne p),
\]
then
\begin{equation}
 \boxed{
  \mu_p(a)=p^2-(25+\varepsilon_a+\delta_a)p+12+\tau_p(a).}
 \label{eq:toric-K3-fiber-count}
\end{equation}
Equivalently, the character sum in
\eqref{eq:toric-fiber-quadratic} is
\begin{equation}
 \boxed{
 \tau_p(a)-(21+\varepsilon_a+\delta_a)p+8.}
 \label{eq:toric-K3-character-sum}
\end{equation}
More precisely, there is a weight-two Frobenius trace $\theta_p(a)$ of
dimension $3-\varepsilon_a-\delta_a$ for which
\begin{equation}
 \boxed{
  \mu_p(a)=p^2-6p+12+\theta_p(a),\qquad
  \mu_p(a)-(p-2)^2=\theta_p(a)-2p+8.}
 \label{eq:toric-K3-deflated-trace}
\end{equation}
Consequently the character sum has absolute value at most $5p+8$, uniformly in
$a\in\mathbb F_p^\times$.
\end{theorem}

\begin{proof}
We give the compactification and all point-count corrections explicitly.
On a nonzero fiber the coordinates used in
Proposition~\ref{prop:toric-mellin-square} identify the original torus fiber
with
\begin{equation}
 \mathcal U_a:\quad
 uv(w-1)(1-uvw)-aw(u-1)(v-1)=0
 \label{eq:toric-K3-affine}
\end{equation}
inside the locus
\[
 u,v,w\ne0,\qquad u,v,w\ne1,\qquad uvw\ne1.
\]
Indeed, the inverse map is
\[
 y=(u-1)^{-1},\qquad z=(v-1)^{-1},\qquad x=-uvw,
\]
and the two remaining exclusions correspond to the factors $1+x$ and
$((1+y)(1+z)+xyz)$ of~$\Lambda$.  Thus
$\#\mathcal U_a(\mathbb F_p)=\mu_p(a)$.

Write $u=U_1/U_0$, $v=V_1/V_0$, and $w=W_1/W_0$.  The closure $X_a$ of
\eqref{eq:toric-K3-affine} in $(\mathbb P^1)^3$ is the multidegree
$(2,2,2)$ hypersurface
\begin{align}
 \mathcal F_a={}&
 U_1V_1(W_1-W_0)(U_0V_0W_0-U_1V_1W_1) \notag\\
 &-aU_0V_0W_0W_1(U_1-U_0)(V_1-V_0)=0.
 \label{eq:toric-K3-homogeneous}
\end{align}

In the shorthand $0,1,\infty$ for points of $\mathbb P^1$, the reduced
boundary $X_a\setminus\mathcal U_a$ consists of the following sixteen
rational curves:
\begin{equation}
\begin{array}{llll}
 u=0,\ v=\infty,&u=0,\ w=\infty,&u=0,\ w=0,&u=0,\ v=1,\\
 v=0,\ u=\infty,&v=0,\ w=\infty,&v=0,\ w=0,&v=0,\ u=1,\\
 u=\infty,\ w=0,&u=\infty,\ w=1,
 &v=\infty,\ w=0,&v=\infty,\ w=1,\\
 u=1,\ w=1,&u=1,\ vw=1,
 &v=1,\ w=1,&v=1,\ uw=1.
\end{array}
\label{eq:toric-K3-boundary-curves}
\end{equation}
Each has $p+1$ rational points.  Their incidence locus contains sixteen
points lying on exactly two components, four points lying on exactly three,
and the single point $(1,1,1)$ lying on four.  There are no other
intersections.  Hence
\begin{equation}
 \#(X_a\setminus\mathcal U_a)(\mathbb F_p)
 =16(p+1)-16-2\cdot4-3=16p-11.
 \label{eq:toric-K3-boundary-count}
\end{equation}

For completeness, we record the full singularity calculation.  The fixed
boundary singularities are
\begin{equation}
\begin{array}{c|c}
\text{points}&\text{type}\\ \hline
 (0,1,\infty),(0,\infty,\infty),
 (1,0,\infty),(\infty,0,\infty)&4A_1,\\
 (0,\infty,0),(\infty,0,0)&2A_2,\\
 (1,1,1)&A_1\ (a\ne1),\quad A_2\ (a=1).
\end{array}
\label{eq:toric-K3-fixed-singularities}
\end{equation}
Here is a local check that also fixes the types in every characteristic
$p\geq5$.  At the first four points, suitable affine coordinates give a
quadratic Hessian of determinant $2a^2$.  At $(0,\infty,0)$ the quadratic
part is
\[
 s(at-r),
\]
and its null direction $(r,s,t)=(az,0,z)$ has cubic coefficient $a^2z^3$;
the other fixed $A_2$ point is symmetric.  At $(1,1,1)$ the quadratic part
is
\[
 -ars-rt-st-t^2,
\]
whose Hessian determinant is $2a(a-1)$.  At $a=1$ it factors as
$-(r+t)(s+t)$, and its null direction $(-z,-z,z)$ has cubic coefficient
$-2z^3$.

It remains to find singularities in the open torus.  Put $q=uvw$ and take
logarithmic derivatives of
\[
 \Lambda=\frac{uv(w-1)(1-uvw)}{w(u-1)(v-1)}.
\]
The three critical equations reduce to
\[
 q(2-u)=q(2-v)=1,
 \qquad qw=1.
\]
They give either the excluded point $(u,v,w)=(1,1,1)$ or
\begin{equation}
 u=v,\qquad w=2-u,\qquad u^2-2u-1=0.
 \label{eq:toric-K3-open-critical-point}
\end{equation}
At the latter point
\[
 a=12u+5,
\]
so such a rational point exists exactly when $a^2-34a+1=0$, and it is
\[
 (u,v,w)=\left(\frac{a-5}{12},\frac{a-5}{12},
                         \frac{29-a}{12}\right).
\]
The logarithmic Hessian determinant at~\eqref{eq:toric-K3-open-critical-point}
is $(u-1)/2$, hence this additional singularity is an $A_1$.  This also
shows directly that the displayed list is complete.

All tangent cones of the $A_2$ points split over $\mathbb F_p$; the
exceptional conic above each rational $A_1$ point has an
$\mathbb F_p$-point and is therefore a rational projective line.  Thus all
the displayed ADE configurations are split over~$\mathbb F_p$.  Resolving
an $A_n$ replaces one rational point by a chain with
\[
 n(p+1)-(n-1)=np+1
\]
rational points, for a net gain of $np$.  If $\widetilde X_a$ denotes the
minimal resolution, then
\begin{equation}
 \#\widetilde X_a(\mathbb F_p)
 =\#X_a(\mathbb F_p)+(9+\varepsilon_a+\delta_a)p.
 \label{eq:toric-K3-resolution-count}
\end{equation}

The hypersurface~\eqref{eq:toric-K3-homogeneous} is anticanonical in
$(\mathbb P^1)^3$ and has only rational double points.  It is normal and
geometrically connected, and the hypersurface exact sequence gives
$H^1(X_a,\mathcal O_{X_a})=0$.  Its ADE resolution is crepant, so
$\widetilde X_a$ is a K3 surface.  Therefore
\[
 \#\widetilde X_a(\mathbb F_p)=1+p^2+\tau_p(a),
 \qquad |\tau_p(a)|\leq22p.
\]
Combining this identity with~\eqref{eq:toric-K3-boundary-count} and
\eqref{eq:toric-K3-resolution-count} proves
\eqref{eq:toric-K3-fiber-count}.  Subtracting the baseline
$(p-2)^2$ in~\eqref{eq:toric-fiber-quadratic} proves
\eqref{eq:toric-K3-character-sum}.

We now use the boundary configuration to deflate all but three of the
generic cohomology classes.  Number the curves in
\eqref{eq:toric-K3-boundary-curves} as $C_1,\ldots,C_{16}$ in the displayed
order, and write $\widetilde C_i$ for their strict transforms.  Each
$\widetilde C_i$ is a smooth rational curve, hence has self-intersection
$-2$.  Away from the rational double points, the nonzero pairwise
intersections are the fourteen transverse pairs
\begin{equation}
\begin{gathered}
 (3,7),\ (2,6),\ (3,4),\ (4,15),\ (1,12),\ (7,8),\ (8,13),\\
 (11,14),\ (12,13),\ (5,10),\ (9,16),\ (10,15),\ (9,11),\ (10,12),
\end{gathered}
\label{eq:toric-K3-smooth-intersection-edges}
\end{equation}
where, for example, $(3,7)$ denotes the edge
$\widetilde C_3\mathbin{-}\widetilde C_7$.

The resolution graph at the four fixed $A_1$ points has exceptional curves
$E_1,\ldots,E_4$ attached respectively to
\begin{equation}
 \{C_2,C_4,C_{16}\},\quad \{C_1,C_2\},\quad
 \{C_6,C_8,C_{14}\},\quad \{C_5,C_6\}.
 \label{eq:toric-K3-A1-attachments}
\end{equation}
At $(0,\infty,0)$ the $A_2$ chain $F_1-F_2$ has $F_1$ attached to
$C_1,C_{11}$ and $F_2$ attached to~$C_3$; at
$(\infty,0,0)$ the chain $G_1-G_2$ has $G_1$ attached to $C_5,C_9$ and
$G_2$ attached to~$C_7$.  For $a\ne1$, the central exceptional curve is
attached to $C_{13},C_{14},C_{15},C_{16}$.  For $a=1$, it is replaced by
an $A_2$ chain $J_1-J_2$, where $J_1$ is attached to $C_{13},C_{16}$ and
$J_2$ to $C_{14},C_{15}$.  The extra $A_1$ exceptional curve when
$a^2-34a+1=0$ is disjoint from this boundary configuration.  These
attachments follow by blowing up the quadratic tangent cones displayed
above.  For instance, at $(0,\infty,0)$ the two exceptional lines are the
two factors of $s(at-r)$; the tangent directions of $C_1,C_{11}$ lie on
$s=0$, whereas that of $C_3$ lies on $at-r=0$.  The other cases are
symmetric, and at $a=1$ the central tangent cone is
$-(r+t)(s+t)$.

This explicit graph already contains a nondegenerate lattice of the needed
rank.  Indeed, put
\[
 \mathcal B=\{\widetilde C_1,\ldots,\widetilde C_{12},
 \widetilde C_{14},\widetilde C_{16},E_1,E_2,E_3,E_4,F_1\}.
\]
Using diagonal entries $-2$ and the edges just listed, its $19\times19$
Gram determinant is~$48$.  If $a=1$, the Gram determinant of
$\mathcal B\cup\{J_1\}$ is $-32$; if
$a^2-34a+1=0$, adjoining the open exceptional curve gives determinant
$-96$.  Thus the divisor classes contain a nondegenerate, Frobenius-stable
subspace of dimension $19+\varepsilon_a+\delta_a$.  Every curve in the
chosen basis is defined over $\mathbb F_p$, so Frobenius acts on this
subspace by~$p$.

Let $\theta_p(a)$ be the Frobenius trace on its orthogonal complement in
$H^2$.  The complement has dimension $3-\varepsilon_a-\delta_a$, and the
Weil bound gives
\[
 \tau_p(a)=(19+\varepsilon_a+\delta_a)p+\theta_p(a),
 \qquad
 |\theta_p(a)|\leq(3-\varepsilon_a-\delta_a)p.
\]
Substitution proves~\eqref{eq:toric-K3-deflated-trace} and the uniform bound
$5p+8$.  Finally, $a=1$ cannot also satisfy
$a^2-34a+1=0$ because the latter value is $-32\ne0$ for $p\geq5$.
\end{proof}

\begin{remark}[Pointwise fiber cancellation is not Mellin
anti-concentration]
\label{rem:toric-K3-Mellin-limit}
Theorem~\ref{thm:toric-fiber-K3} improves the elementary termwise Weil
bound for~\eqref{eq:toric-fiber-quadratic} from $O(p^{3/2})$ to the explicit
bound $5p+8$.
It does not bound the number of vanishing Ap\'ery coefficients.  Indeed,
for $1\leq j\leq p-2$,
\[
 b_j=-\sum_{a\ne0}\mu_p(a)a^j
\]
is the defining-characteristic multiplicative Mellin coefficient of the
entire fiber-count function.  The constant baseline has zero nontrivial
Mellin transform, while an $O(p)$ pointwise bound on the remaining trace
function does not prevent exact cancellation modulo~$p$.  The example
$(p,j)=(11,5)$ already gives such a cancellation.  Thus the K3 collapse is
a genuine pointwise theorem, but the missing input remains crystalline
Mellin anti-concentration.
\end{remark}
